\documentclass[11pt,a4paper]{article}

\usepackage[a4paper,margin=1in]{geometry}
\usepackage{fontspec}
\IfFontExistsTF{Times New Roman}{\setmainfont{Times New Roman}}{\setmainfont{TeX Gyre Termes}}
\IfFontExistsTF{Menlo}{\setmonofont{Menlo}}{\setmonofont{Latin Modern Mono}}
\usepackage{amsmath,amssymb}
\usepackage{booktabs}
\usepackage{longtable}
\usepackage{array}
\usepackage{enumitem}
\usepackage{graphicx}
\usepackage{fancyhdr}
\usepackage[round,authoryear]{natbib}
\usepackage[hidelinks]{hyperref}
\usepackage{xurl}
\usepackage{microtype}

\setlength{\parindent}{0pt}
\setlength{\parskip}{0.65em}
\setlist[itemize]{leftmargin=1.4em}
\newcommand{\file}[1]{\path{#1}}
\newcommand{\code}[1]{\path{#1}}
\newcolumntype{L}[1]{>{\raggedright\arraybackslash}p{#1}}

\pagestyle{fancy}
\fancyhf{}
\lhead{\small v0 Dry-Run Dataset}
\rhead{\small Dry-Run Review}
\cfoot{\thepage}
\renewcommand{\headrulewidth}{0.4pt}

\title{\textbf{v0 Dry-Run Dataset Report}\\
\large Eight-Regime Surrogate Dataset Generation Test}
\author{Codex-assisted surrogate data review}
\date{June 18, 2026}

\begin{document}
\maketitle

\begin{abstract}
This report reviews the v0 dry-run dataset generation test for the American option risk-surface
project. The dry run generated a small eight-regime package using the validated baseline
American Crank--Nicolson / projected SOR solver. It produced one compressed \code{.npz} file,
audit manifests, diagnostics, and optional review figures. It did not generate the full
288-regime dataset, train a surrogate model, create \code{.pt} files, or approve production
Greek labels. The dry-run review decision is
\code{READY_FOR_V1_SMALL_GRID_PLANNING}.
\end{abstract}

\tableofcontents
\newpage

\section{Purpose of v0 Dry Run}

The v0 dry run is the first controlled execution after the surrogate dataset construction plan.
Its purpose is to test the proposed pipeline on eight regimes before any larger generation step.
The dry run checks whether solver outputs, continuation premium, boundary metadata, Greek masks,
LCP diagnostics, split metadata, and file manifests can be assembled into one auditable package.

This is a dry-run generation test only. It is not the final dataset, not a neural-surrogate
training stage, and not approval for broad production-scale sweeps.

\section{Roadmap and Construction-Plan Link}

The compressed roadmap places this work in Stage 1: v0 dry-run dataset plus dry-run review. The
preceding dataset construction plan named the full 288-regime grid, split policy, failure rules,
and eight-regime dry-run subset. This report records the outcome of executing that subset.

The solver source remains the validated classical benchmark:
\begin{center}
\code{solver_name = american_crank_nicolson_psor_price}\\
\code{solver_variant = baseline_cn_psor}
\end{center}
The numerical method and American-option setting follow the project validation ladder and the
classical option-pricing references already integrated into the solver reports
\citep{black_scholes_1973,merton_1973,brennan_schwartz_1977,wilmott_howison_dewynne_1995}.

\section{Exact Eight Regimes}

{\scriptsize
\setlength{\tabcolsep}{1.5pt}
\begin{longtable}{L{0.07\textwidth} L{0.41\textwidth} L{0.06\textwidth} L{0.05\textwidth} L{0.06\textwidth} L{0.05\textwidth} L{0.05\textwidth} L{0.06\textwidth}}
\toprule
Dry ID & Regime ID & Type & \(T\) & \(\sigma\) & \(r\) & \(q\) & Split \\
\midrule
\endhead
\code{dry_01} & \code{put_T100_s020_r005_q003} & put & 1.0 & 0.20 & 0.05 & 0.03 & train \\
\code{dry_02} & \code{put_T100_s060_r005_q003} & put & 1.0 & 0.60 & 0.05 & 0.03 & train \\
\code{dry_03} & \code{put_T100_s020_r001_q003} & put & 1.0 & 0.20 & 0.01 & 0.03 & train \\
\code{dry_04} & \code{put_T200_s040_r010_q010} & put & 2.0 & 0.40 & 0.10 & 0.10 & test \\
\code{dry_05} & \code{call_T100_s020_r005_q006} & call & 1.0 & 0.20 & 0.05 & 0.06 & train \\
\code{dry_06} & \code{call_T100_s020_r005_q010} & call & 1.0 & 0.20 & 0.05 & 0.10 & train \\
\code{dry_07} & \code{call_T100_s060_r005_q006} & call & 1.0 & 0.60 & 0.05 & 0.06 & train \\
\code{dry_08} & \code{call_T100_s020_r005_q000} & call & 1.0 & 0.20 & 0.05 & 0.00 & test \\
\bottomrule
\end{longtable}
}

All regimes used \(K=1\), \(S_{\max}=4\), and \(M=N=120\). The no-dividend call row is a
control regime, not dividend-call early-exercise evidence.

\section{Feature, Label, and Mask Schema}

The stored feature matrix uses:
\[
x=\left(\log(S/K),\ \tau/T,\ r,\ q,\ \sigma,\ T,\ \mathrm{is\_call}\right).
\]
The package stores reliable price-style labels and diagnostic labels separately:
\begin{itemize}
  \item core dry-run labels: \code{value_over_K}, \code{payoff_over_K}, and \code{premium_over_K};
  \item diagnostic labels: \code{exercise_indicator} and \code{boundary_spot_over_K};
  \item cautioned Greek diagnostics: \code{delta} and \code{scaled_gamma}.
\end{itemize}

Masks are included for payoff-kink-near, boundary-near, maturity-row, strict-interior,
Gamma-allowed, Delta-allowed, exercise-region, and continuation-region nodes. These masks are
mandatory metadata. They are especially important because Gamma remains fragile near kinks,
maturity rows, and extracted boundary bands.

\section{Diagnostics and Acceptance Checks}

A regime is accepted into the dry-run arrays only when PSOR converges, obstacle violation is at
most \(10^{-8}\), equation violation is at most \(10^{-6}\), complementarity product is at most
\(10^{-6}\), and required metadata and masks are present. Failed regimes would remain visible in
the CSV manifests with \code{review_required}; they would not be silently included as accepted
label rows.

\begin{longtable}{L{0.36\textwidth} L{0.22\textwidth} L{0.30\textwidth}}
\toprule
Metric & Dry-run result & Interpretation \\
\midrule
\endhead
Regimes solved & 8 & Exactly the approved dry-run subset. \\
Accepted regimes & 8 & All regimes passed the dry-run acceptance checks. \\
Accepted sampled rows & 41,624 & 5,203 rows per accepted regime. \\
Train rows & 31,218 & Six train regimes. \\
Test rows & 10,406 & Two test/control regimes. \\
Maximum obstacle violation & \(0.0\) & Payoff obstacle respected in accepted regimes. \\
Maximum equation violation & \(4.90\times 10^{-8}\) & Below the \(10^{-6}\) diagnostic threshold. \\
Maximum complementarity product & \(5.34\times 10^{-10}\) & Below the \(10^{-6}\) diagnostic threshold. \\
Maximum PSOR iterations & 41 & Highest effort occurs in a high-volatility put regime. \\
Total solver runtime & about 0.82 seconds & Small dry run is computationally manageable. \\
\bottomrule
\end{longtable}

\section{Output Files and Manifests}

\begin{longtable}{L{0.38\textwidth} L{0.14\textwidth} L{0.34\textwidth}}
\toprule
Output & Type & Purpose \\
\midrule
\endhead
\file{results/04_surrogate_dataset/v0_dry_run/dataset_v0_dry_run.npz} & \code{.npz} & Accepted-row dry-run arrays only. \\
\file{results/04_surrogate_dataset/v0_dry_run/regime_manifest.csv} & CSV & Regime parameters, convergence, and acceptance status. \\
\file{results/04_surrogate_dataset/v0_dry_run/diagnostic_summary.csv} & CSV & PSOR, LCP, boundary, Greek, and runtime diagnostics. \\
\file{results/04_surrogate_dataset/v0_dry_run/split_assignment.csv} & CSV & Regime-level split assignment and row counts. \\
\file{results/04_surrogate_dataset/v0_dry_run/schema_snapshot.csv} & CSV & Feature, label, mask, and audit-field schema. \\
\file{results/04_surrogate_dataset/v0_dry_run/output_manifest.csv} & CSV & Output audit manifest for all dry-run artifacts. \\
\bottomrule
\end{longtable}

The compressed package contains the keys \code{X}, \code{y_value}, \code{y_payoff},
\code{y_premium}, \code{y_exercise_indicator}, \code{y_boundary}, \code{y_delta},
\code{y_scaled_gamma}, \code{masks}, \code{regime_index}, \code{feature_names},
\code{label_names}, \code{mask_names}, \code{audit_numeric}, \code{audit_numeric_names},
\code{regime_ids}, \code{dry_run_ids}, and \code{split_names}.

\section{Figure Interpretation}

Three optional figures were created for human review. They are diagnostic figures only, not
stress-map datasets.

\begin{figure}[htbp]
\centering
\includegraphics[width=0.82\textwidth]{results/04_surrogate_dataset/v0_dry_run/figures/v0_dry_run_row_counts.png}
\caption{Accepted sampled rows by dry-run regime. Each accepted regime contributes 5,203 sampled rows in the reporting region.}
\end{figure}

\begin{figure}[htbp]
\centering
\includegraphics[width=0.82\textwidth]{results/04_surrogate_dataset/v0_dry_run/figures/v0_dry_run_diagnostics.png}
\caption{Obstacle, equation, and complementarity diagnostics by regime. The equation residuals remain below the dry-run threshold.}
\end{figure}

\begin{figure}[htbp]
\centering
\includegraphics[width=0.78\textwidth]{results/04_surrogate_dataset/v0_dry_run/figures/v0_dry_run_premium_distribution.png}
\caption{Distribution of continuation premium over accepted rows. This checks the dry-run label pipeline without creating broad surfaces.}
\end{figure}

\section{Dry-Run Review Decision}

The dry-run review decision is:
\begin{center}
\fbox{\code{READY_FOR_V1_SMALL_GRID_PLANNING}}
\end{center}

This means the eight-regime pipeline is ready for a v1 small-grid planning step. It does not mean
the full dataset is approved for generation, and it does not approve surrogate training.

\section{Limitations}

\begin{itemize}
  \item The package covers eight regimes only, not the planned 288-regime grid.
  \item Boundary labels remain threshold-based diagnostics, not exact analytical free boundaries.
  \item Delta and scaled Gamma are diagnostic fields. Gamma remains high caution and should be
  used only with masks.
  \item Runtime observations from eight regimes do not prove large-scale generation cost.
  \item The no-dividend call control has no boundary found, which is expected control behavior
  rather than a failed dividend-call boundary.
\end{itemize}

\section{What This Dry Run Supports}

The dry run supports the claim that the approved eight-regime generation pipeline can assemble
features, reliable price and premium labels, diagnostic labels, masks, solver diagnostics, split
metadata, and manifests into one auditable package. It also supports moving to a v1 small-grid
planning discussion.

\section{What This Dry Run Does Not Support}

The dry run does not support final dataset claims, neural surrogate training, production Greek
labels, broad parameter-sweep conclusions, or final paper claims about all regimes. It is a
pipeline test and review gate.

\section{Recommended Next Step}

The recommended next step is a v1 small-grid planning gate. That gate should decide whether to
generate all 288 planned regimes, whether any additional higher-grid confirmation is required,
and whether Delta or Gamma should remain metadata-only in the first generated version.

\bibliographystyle{plainnat}
\bibliography{reports/03_solver/references}

\end{document}
